\documentclass[10pt,a4paper]{amsart}
\usepackage[margin=19mm]{geometry}
\usepackage{amsmath,amssymb,mathtools}
\usepackage[scr=rsfs,cal=euler]{mathalfa}
\usepackage{xfrac}
\usepackage[unicode,hidelinks,bookmarks=false]{hyperref}
\newcommand{\ie}{i.e.\ }
\newcommand{\cf}{cf.\ }
\newcommand{\resp}{respectively\ }
\newcommand{\Subsection}{\subsection}
\providecommand{\nobreakdash}{\nobreak\hbox{-}}
\setlength{\emergencystretch}{2em}
\multlinegap0pt
\usepackage{amssymb}
\usepackage[shortlabels]{enumitem}
\newtheorem{lemma}{Lemma}
\newtheorem{proposition}{Proposition}
\newcommand{\FCC}{L_{\mathrm F}}
\newcommand{\Id}{I}
\newcommand{\PP}{\mathcal P_3^1}
\newcommand{\R}{{\mathbb R}}
\newcommand{\Th}{\Theta}
\newcommand{\frob}{\mathrm F}
\newcommand{\op}{\mathrm{op}}
\newcommand{\tr}{\operatorname{tr}}
\title{On Sarnak--Str\"{o}mbergsson conjecture}
\author{Senping Luo, Juncheng Wei}
\date{Source: arXiv:2609.17356v1 (CC BY 4.0)}
\begin{document}
\maketitle

\noindent\emph{Source and licence.} On Sarnak--Str\"{o}mbergsson conjecture. Version arXiv:2609.17356v1 (CC BY 4.0), distributed by arXiv under the Creative Commons Attribution 4.0 International licence, http://creativecommons.org/licenses/by/4.0/, as recorded in the arXiv licence metadata for this exact version and shown on the abstract page https://arxiv.org/abs/2609.17356v1. Full licence notice: \texttt{licenses/arxiv-2609.17356-CC-BY-4.0.txt}.\par\smallskip

\noindent\textit{Editorial scope.} Counted unit: the article's proposition T:prop:smallballlarge (source0.tex lines 1433--1435). The article's complete original proof of it is delivered with it (source0.tex lines 1436--1511, one \texttt{proof} environment of 76 lines). No mathematics is written, repaired or re-derived: the statement and the proof are the article's own lines, in the article's own notation, and no retained line is substituted or re-flowed.  Retained uncounted as the source-local context the proof needs: lines 902--904; lines 1224--1236; lines 1333--1336; lines 1388--1398.  Nothing was omitted from any retained span, so the package carries no omission record.  No retained line contains a citation, so the wrapper carries no bibliography and no bibliography entry.  No retained line contains a figure environment, an image-inclusion command or drawing code, so no image is delivered and the package declares no asset dependency.  Editorial formatting only, in addition: an AMS \texttt{amsart} preamble that restates the article's own declarations and macros used by the retained lines, namely the environments \texttt{lemma}, \texttt{proposition} on separate counters, together with the standard packages those lines need.  The retained environments print in this document as Lemma 1, Proposition 1 in order of appearance, whereas the article prints the same items with the numbering of its own full document.  The complete native source of the article as distributed in the arXiv e-print for this exact version is bundled unchanged as \texttt{sources/210590/source0.tex}.
\begin{equation}\label{T:eq:isotropic}
 \sum_{|v|^2=n\rho}vv^T=\frac{n\rho N_n}{3}\Id.
\end{equation}

\begin{lemma}[Trace control on the determinant-one surface]\label{T:lem:det}
Suppose $A=\Id+T\in\PP$ and $\|T\|_{\op}\le\delta<1$, with $\delta\ge0$. Then
\begin{equation}\label{T:eq:tracebounds}
 \frac{\|T\|_{\frob}^2}{2(1+\delta)}
 \le\tr T\le
 \frac{\|T\|_{\frob}^2}{2(1-\delta)}.
\end{equation}
Put $\eta=\tr T/3$ and $S=T-\eta\Id$. If $\delta<2/3$, then
\begin{equation}\label{T:eq:etabounds}
 0\le\eta\le\frac{\|S\|_{\frob}^2}{6-9\delta},
 \qquad \|S\|_{\frob}\le\sqrt3\,\delta.
\end{equation}
\end{lemma}

\begin{equation}\label{T:eq:localquant}
 \Th(\alpha,A^{1/2}\FCC)-\Th(\alpha,\FCC)
 \ge\frac{u}{50}e^{-u}\|A-\Id\|_{\frob}^2.
\end{equation}

\begin{lemma}[Six-point moment bound]\label{T:lem:mgf}
Let $x_1,\ldots,x_6\in\R$ have mean zero. Define $\sigma\ge0$ and $M$ by
\[
 \sigma^2=\frac16\sum_{j=1}^6x_j^2,
 \qquad M=\frac16\sum_{j=1}^6e^{x_j}.
\]
Then
\begin{equation}\label{T:eq:mgf}
 \log M\ge\frac1{13}\min(\sigma^2,\sigma).
\end{equation}
\end{lemma}

\begin{proposition}[A one-percent FCC ball, unbounded scales]\label{T:prop:smallballlarge}
Let $A\in\PP$ and $u=\pi\rho\alpha$. Inequality \eqref{T:eq:localquant} holds when $u\ge64$ and $\|A-\Id\|_{\op}\le1/100$.
\end{proposition}

\begin{proof}
We separate the scalar and trace-free parts of the strain. Write
$A-\Id=T=\eta\Id+S$ as in Lemma~\ref{T:lem:det}. Then
\begin{equation}\label{T:eq:eta1percent}
 \eta\le\frac15\|S\|_{\frob}^2,
 \qquad \|S\|_{\frob}<\frac1{50}.
\end{equation}
The contribution of the first shell takes the exact form
\begin{equation}\label{T:eq:firstmgf}
 12e^{-u}e^{-u\eta}M,
 \qquad
 M=\frac16\sum_{k,\pm}
 \exp\left(\frac u2S_{kk}\pm uS_{ij}\right),
\end{equation}
where $k$ runs through $1,2,3$. For each $k$, the indices $i<j$
are the two complementary indices, so that
\[
 (k,i,j)=(1,2,3),\quad (2,1,3),\quad (3,1,2).
\]
The notation $\sum_{k,\pm}$ means that both signs are included for each
of these three triples. Thus the sum contains exactly six terms;
there is no additional independent summation over $i$ or $j$.
Explicitly,
\[
\begin{aligned}
 M=\frac16\Bigl[{}
 &e^{uS_{11}/2+uS_{23}}+e^{uS_{11}/2-uS_{23}}\\
 &+e^{uS_{22}/2+uS_{13}}+e^{uS_{22}/2-uS_{13}}\\
 &+e^{uS_{33}/2+uS_{12}}+e^{uS_{33}/2-uS_{12}}
 \Bigr].
\end{aligned}
\]
Thus $M$ is the mean of the six displayed exponentials. Their
exponents sum to zero: $\tr S=0$, and the two off-diagonal terms
cancel for each $k$. The corresponding variance is
\begin{equation}\label{T:eq:variance}
 \sigma^2=\frac{u^2}{12}\bigl(D(S)+4O(S)\bigr)
 \ge\frac{u^2}{12}\|S\|_{\frob}^2.
\end{equation}
The key estimate is $\log M\ge2u\eta$. We prove it in two cases.
If $\sigma\le1$, Lemma~\ref{T:lem:mgf} gives
\[
 \log M\ge\frac{u^2}{156}\|S\|_{\frob}^2
 \ge\frac25u\|S\|_{\frob}^2\ge2u\eta,
\]
since $64/156>2/5$. If $\sigma\ge1$, then
\[
 \log M\ge\frac{u\|S\|_{\frob}}{13\sqrt{12}}
 >\frac{u\|S\|_{\frob}}{52}
 \ge\frac25u\|S\|_{\frob}^2\ge2u\eta,
\]
where \eqref{T:eq:eta1percent} was used in the penultimate inequality.

By \eqref{T:eq:firstmgf}, this proves that the first-shell increase
is at least $12e^{-u}(e^{u\eta}-1)\ge12u\eta e^{-u}$.
For the higher shells, the linear convexity bound and
\eqref{T:eq:isotropic} show that the total possible decrease is at
most
\[
 u\eta e^{-u}\sum_{n\ge2}nN_nq^{n-1}.
\]
The coefficient satisfies
\begin{equation}\label{T:eq:highertail3}
 \sum_{n\ge2}nN_nq^{n-1}
 \le30\left(\frac{1+q}{(1-q)^3}-1\right)<3
 \qquad(q\le1/49).
\end{equation}
Subtracting this possible loss leaves an increase of at least
$9u\eta e^{-u}$. Finally, the lower trace estimate gives
\[
 \eta\ge\frac{\|T\|_{\frob}^2}{6(1+1/100)}
 \ge\frac17\|T\|_{\frob}^2,
\]
which yields a bound stronger than \eqref{T:eq:localquant} and
completes the proof.
\end{proof}

\end{document}
